\documentclass[11pt,a4paper]{article}


\usepackage{cite}
\usepackage{amsmath,amssymb,amsfonts}
\usepackage{amsthm}
\usepackage{graphicx}
\usepackage{algorithm}
\usepackage[hidelinks]{hyperref}
\usepackage{textcomp}

\usepackage[utf8]{inputenc} % allow utf-8 input
\usepackage[T1]{fontenc}    % use 8-bit T1 fonts
\usepackage{amsmath,amssymb,amsfonts}
\usepackage{algpseudocode}
\usepackage{todonotes}
\usepackage{multirow}
\usepackage{url}
\usepackage{cite}
\usepackage{textcomp}
\usepackage{xcolor}
\usepackage{optidef}
\usepackage{arydshln}
\usepackage{balance}
\usepackage{tikz}
\usepackage{circuitikz}
\usetikzlibrary{arrows.meta}
\tikzstyle{block} = [draw, rectangle, 
    minimum height=3em, minimum width=3em]
\tikzstyle{sum} = [draw, circle, node distance=1cm]
\tikzstyle{input} = [coordinate]
\tikzstyle{output} = [coordinate]
\tikzstyle{pinstyle} = [pin edge={to-,thin,black}]

\graphicspath{{./figures/}}



\newtheorem{theorem}{Theorem}
\newtheorem{lemma}[theorem]{Lemma}
\newtheorem{corollary}{Corollary}[theorem]
\newtheorem{remark}{Remark}
\newtheorem{proposition}{Proposition}[section]
%
\newcommand{\dsp}{\displaystyle}
\newcommand{\BBM}{\begin{bmatrix}}
\newcommand{\EBM}{\end{bmatrix}}
\newcommand{\BEQ}{\begin{equation}}
\newcommand{\EEQ}{\end{equation}}
\newcommand{\BIT}{\begin{itemize}}
\newcommand{\EIT}{\end{itemize}}
\newcommand{\cf}{{\it cf.}}
\newcommand{\eg}{{\it e.g.}}
\newcommand{\ie}{{\it i.e.}}
\newcommand{\st}{\mbox{subject to}}
\newcommand{\ones}{\mathbbm{1}}
\newcommand{\reals}{\mathbb{R}}
\newcommand{\ereals}{\bar{\reals}}
\newcommand{\pereals}{(-\infty,+\infty]}
\newcommand{\naturals}{\mathbb{N}}
\newcommand{\symm}{\mathbb{S}}
\newcommand{\quadcone}{\mathbb{Q}}
\newcommand{\expcone}{K_{\exp}}
\newcommand{\complex}{\mathbb{C}}
\newcommand{\integers}{\mathbb{Z}}
\newcommand{\psimplex}{\Delta}
\newcommand{\eqbydef}{\mathrel{\stackrel{\Delta}{=}}}
\newcommand{\im}{\mathrm{i}}
\newcommand{\minimize}{\mbox{minimize}}
\newcommand{\maximize}{\mbox{maximize}}
\newcommand{\Range}{\mathcal{R}}
\newcommand{\nullspace}{\mathcal{N}}
\newcommand{\feasset}{\mathcal{F}}
\newcommand{\probdom}{\mathcal{D}}

\newcommand{\argit}[2]{#1_{#2}}
\newcommand{\xit}[1]{\argit{x}{#1}}
\newcommand{\majf}{\bar f}
\newcommand{\maj}[2]{\majf(#1\,\vert\,#2)}
\newcommand{\indicator}[1]{I_{#1}}


\DeclareMathOperator{\Span}{span}
\DeclareMathOperator{\nullity}{nullity}
\DeclareMathOperator{\rank}{rank}
\DeclareMathOperator{\Tr}{tr}
\DeclareMathOperator{\diag}{diag}
\DeclareMathOperator{\blkdiag}{blkdiag}
\DeclareMathOperator{\vect}{vec}
\DeclareMathOperator{\svect}{svec}
\DeclareMathOperator{\dom}{dom}
\DeclareMathOperator{\conv}{conv}
\DeclareMathOperator{\relint}{relint}
\DeclareMathOperator{\interior}{int}
\DeclareMathOperator{\aff}{aff}
\DeclareMathOperator{\cone}{cone}
\DeclareMathOperator{\epi}{epi}
\DeclareMathOperator{\cl}{cl}
\DeclareMathOperator*{\argmin}{argmin}
\DeclareMathOperator*{\argmax}{argmax}
\DeclareMathOperator{\sgn}{sgn}
\DeclareMathOperator{\prox}{prox}
\DeclareMathOperator{\nnz}{nnz}

\DeclareMathOperator{\sqr}{sqr}
\DeclareMathOperator{\adj}{adj}
\DeclareMathOperator{\bdiag}{bdiag}

\DeclareMathOperator{\Expect}{\mathbb{E}}
\DeclareMathOperator{\Var}{Var}

\newcommand{\E}[2][]{\Expect_{#1}\!\left[\,#2\,\right]}   % Expectation with argument
\newcommand{\CE}[2]{\E{#1 \,\middle|\, #2}}                              % Conditional expectation

\DeclareMathOperator{\Probability}{\mathbb{P}}
\newcommand{\Prob}[2][]{\Probability_{#1}\!\left[\,#2\,\right]}     % Probability with argument
\newcommand{\CProb}[2]{\Prob{#1 \,\middle|\, #2}}                       % Conditional probability

\newcommand{\tnabla}{\widetilde{\nabla}}


\newcommand{\mfile}[1]{\txtmint{#1}}
\newcommand{\mfilelabel}[2]{\txtmint{#1}}



\title{Note on informativity of $y_b$}

\begin{document}
\maketitle

\section{Augmented $r$}

\todo[inline]{I will rewrite it to use $\tilde y^B$ instead of $\bar y^B$.}

Defining $\bar y^B$ as the subvector of $y^B$ containing only the outputs that connect to $A$, and let $\bar \Upsilon_{AB}$ be the corresponding submatrix of $\Upsilon_{AB}$ with zero rows removed, this allow us to define $\bar \Omega = \BBM \Omega_A & \bar\Upsilon_{AB} \EBM$, and
\begin{align}\label{eq:aug_r_A}
    \bar r = \BBM r^A \\ \bar y^{B} \EBM.
\end{align}

\todo[inline]{Note that we don't need $y^B$ to be informative, only the subset of $y^B$ that is connected to $A$, which we call $\bar y^B$ here.}


\section{Conditions for $r^B$ for informative experiments}
Here we will derive sufficient conditions for $r^B$ to ensure informative experiments, we will assume $\Phi_{r^A} \succ 0$. Let $|\bar B|$ denote the number of outputs in $B$ that connect to $A$.

\subsection{Frequency–Domain Rank Condition}

We start by characterizing when $\bar y^{B}$ provides sufficient excitation to complement $r^A$.

\begin{lemma}\label{lem:informative_rank_lemma}Let $\mathbf u(z) \in \mathbb{C}^m$ be a quasi-stationary signal with power spectral density $\Phi_{u}(\omega) \succ 0$ for almost all $\omega \in [-\pi, \pi]$. Let $\mathbf F(z) \in \mathbb{C}^{p \times m}$ be a proper rational transfer function with all poles inside the
unit circle, and consider the signal $\mathbf y(z) = \mathbf F(z)\mathbf u(z)$.  Then the spectrum of $y$ satisfies $\Phi_{y}(\omega) \succ 0$ for almost all $\omega$, if and only if $\mathbf F(z)$ has full row rank for almost all $z$.
\end{lemma}

\begin{proof}
The power spectral density of $y$ is given by
\begin{align*}
    \Phi_{y}(\omega) = \mathbf F(e^{j\omega})\Phi_{u}(\omega)\mathbf F^*(e^{j\omega}),
\end{align*}
where $\mathbf F^*(e^{j\omega}) = \mathbf F^T(e^{-j\omega})$ denotes the conjugate transpose. 

Suppose $\mathbf F(z)$ has full row rank $p$ for almost all $z$. Then for any nonzero $\mathbf v \in \mathbb{R}^p$, the row vector $\mathbf v^T\mathbf F(z)$ is not identically zero as a rational function. Since rational functions have at most a finite number of zeros, $\mathbf v^T\mathbf F(e^{j\omega}) \neq 0$ for almost all $\omega \in [-\pi, \pi]$. Therefore,
\begin{align*}
    \mathbf v^T\Phi_{\mathbf y}(\omega)\mathbf v = \|\mathbf F^*(e^{j\omega})\mathbf v\|_{\Phi_{\mathbf u}(\omega)}^2 > 0
\end{align*}
for almost all $\omega$, where the inequality follows from $\mathbf F^*(e^{j\omega})\mathbf v \neq 0$ and $\Phi_{\mathbf u}(\omega) \succ 0$. Since this holds for all nonzero $\mathbf v$, we have $\Phi_{\mathbf y}(\omega) \succ 0$ for almost all $\omega$.


Suppose $\Phi_{\mathbf y}(\omega) \succ 0$ for almost all $\omega$. Assume for contradiction that $\mathbf F(z)$ does not have full row rank. Then there exists a nonzero $\mathbf v \in \mathbb{R}^p$ such that $\mathbf v^T\mathbf F(z) = 0$ for all $z$. In particular, $\mathbf v^T\mathbf F(e^{j\omega}) = 0$ for almost all $\omega$, which implies
\begin{align*}
    \mathbf v^T\Phi_{\mathbf y}(\omega)\mathbf v = \|\mathbf F^*(e^{j\omega})\mathbf v\|_{\Phi_{\mathbf u}(\omega)}^2 = 0
\end{align*}
for almost all $\omega$, contradicting $\Phi_{\mathbf y}(\omega) \succ 0$.
\end{proof}

Define $\mathbf{ T_{B}}(z)$ as the transfer function
\begin{align*}
    \mathbf{ T_{B}}(z) = \BBM \mathbf{T_{BA}}(z) & \mathbf{T_{BB}}(z) \EBM,
\end{align*}
where $\mathbf{T_{BA}}$ and $\mathbf{T_{BB}}$ are respectively the transfer functions from $r^A$ and $r^B$ to $\bar y^B$. 

Using Lemma~\ref{lem:informative_rank_lemma}, the signal $\bar y^B$ has a positive definite spectrum if 
\begin{align}\label{eq:rank-condition}
    \text{rank}\bigl(\mathbf{ T_{B}}(e^{j\omega})\bigr)=|\bar B| \quad\text{for a.e.\ }\omega\in[-\pi,\pi].
\end{align}

Condition~\eqref{eq:rank-condition} is the \emph{frequency-domain rank test} for the informativity of the outputs of $B$ that are connected to $A$.

The following theorem provides sufficient excitation conditions for $r^B$ in order to $\Phi_{y^B}(\omega) \succ 0$.


\begin{theorem}\label{thm:minimal_excitation}
Given that $\Phi_{r^A}(\omega) \succ 0$ for almost all $\omega$, the outputs of $B$ connected to $A$ satisfies the condition $\Phi_{\bar y^B}(\omega) \succ 0$ if there exists a subset $\bar r^B \subseteq r^B$ of the external signals applied to $B$ such that:
\begin{enumerate}
    \item \textbf{Rank Condition:} The transfer function
    \begin{align}\label{eq:rank_condition}
        \mathbf{T_{\bar B}}(z) = \BBM \mathbf{T_{BA}}(z) & \mathbf{T_{B\bar B}}(z) \EBM,
    \end{align}
     has full row rank $|\bar B|$ for almost all $\omega \in [-\pi, \pi]$, where $\mathbf{T_{B\bar B}}(z)$ denotes the transfer function from the selected external signals in $\bar r^B$ to $\bar y^B$.
    
    \item \textbf{Excitation Condition:} The selected external signals $\bar r^B$ satisfy
    \begin{align}\label{eq:excitation_condition}
        \Phi_{\bar r^B}(\omega) \succ 0 \quad \text{for almost all } \omega \in [-\pi, \pi].
    \end{align}
\end{enumerate}
\end{theorem}

\begin{proof}
    We can write
\begin{align*}
    \bar y^B &= \mathbf{T_{BA}}(z) \mathbf r^A + \mathbf{T_{BB}}(z) \mathbf r^B + \text{noise terms}.
\end{align*}

For the signal $\bar y^B$ to satisfy $\Phi_{\bar y^B}(\omega) \succ 0$, $ r^B$ must provides sufficient excitation to complement $r^A$, i.e. that some subset $\bar r^B$ of external signals provides excitation to $\bar y^B$ through paths that are linearly independent from the paths taken by $r_A$.

Let $\bar r^B$ be a subset of $r^B$ that satisfies the rank condition. Then
\begin{align*}
   \bar y^B &= \mathbf{T_{BA}}(z) \mathbf r^A + \mathbf{T_{B\bar B}}(z) \mathbf{\bar r}^B + \mathbf{T_{B\tilde B}}(z) \mathbf{\tilde r}^{B} + \text{noise terms},
\end{align*}
where $\tilde r^{ B}$ represents the set of remaining external signal in $r^B$ that are not in $\bar r ^B$, i.e., $\mathbf{T_{BB}}(z) \mathbf r^B = \mathbf{T_{B\bar B}}(z) \mathbf{\bar r}^B + \mathbf{T_{B\tilde B}}(z) \mathbf{\tilde r}^{B}$.

By the rank condition, the matrix $\mathbf{T_{\bar B}}(z)$ in \eqref{eq:rank_condition} has full row rank, which by Lemma~\ref{lem:informative_rank_lemma} ensures that the component $\mathbf{T_{BA}}(z) \mathbf r^A + \mathbf{T_{B\bar B}}(z) \mathbf{\bar r}^B$ span the full output space of $\bar y^B$.

The excitation condition \eqref{eq:excitation_condition}, combined with the assumption that $\Phi_{r^A}(\omega) \succ 0$, ensures that $\bar r^B$ have sufficient spectral richness to guarantee $\Phi_{\bar y^B}(\omega) \succ 0$.
\end{proof}


\begin{corollary}\label{cor:minimal_subset}
The minimal number of external signals in $\bar r^B$ required to satisfy the rank condition in Theorem~\ref{thm:minimal_excitation} is
\begin{align*}
    |\bar r^B|_{\min} = \max\{0, |\bar B| - \text{rank}(\mathbf{T_{BA}}(e^{j\omega}))\}
\end{align*}
for almost all $\omega$.
\end{corollary}



\subsection{Path-Based Verification}
Verifying the rank condition in Theorem~\ref{thm:minimal_excitation} directly is intractable when the transfer functions are unknown. Instead, we need a purely graph-theoretic condition. The link between the generic rank of transfer functions and vertex-disjoint paths in the underlying graph was established in \cite{graph_TF_rank}. Two or more paths in a directed graph are said to be \emph{vertex-disjoint} if they do not share any vertex.


\begin{theorem}\label{th:graph_TF_rank}
Consider a Multiple-Input Multiple-Output system
\begin{align*}
    \xi_{k+1} &= F\xi_{k} + Gu_k\\
    y_k &= H\xi_{k},
\end{align*}
with output $y_k\in \reals^p$, input $u_k\in \reals^m$, and state $\xi_k\in \reals^n$. The underlying graph of the system is defined by the set of vertices $V = U\cup X \cup Y$, where
\begin{align*}
    U &= \{u^1, \dots, u^m\}, \\ 
    X &= \{\xi^1, \dots, \xi^n\},\\
    Y &= \{y^1, \dots, y^n\},
\end{align*}
and the set of edges
\begin{align*}
    E &= \{(u^j,\xi^i)|G_{ij}\neq 0\} \cup \{(\xi^j,\xi^i)|F_{ij}\neq 0\}\\ &\quad \cup \{(\xi^j,y^i)|H_{ij}\neq 0\}.
\end{align*}
Then, the generic rank of the transfer function $\mathbf{G}(z) = H(zI - F)^{-1}G \in \mathbb{C}^{m\times p}$ is equal to the maximum number of vertex-disjoint paths from the input nodes $U$ to the output nodes $Y$ in the underlying graph of the system.
\end{theorem}
\begin{proof}
   See \cite[Theorem 2]{graph_TF_rank}. 
\end{proof}

The following corollary follows directly from the theorem.

\begin{corollary}
    The maximum number of vertex disjoint paths from the external inputs in $ r^A$ and $\bar r^B$ to $\bar y^B$ is equal to the generic rank of $\mathbf{T_{\bar B}}(e^{j\omega})$.
\end{corollary}

The rank condition in Theorem~\ref{thm:minimal_excitation} can then be verified generically using graph-theoretic conditions.

\begin{proposition}\label{prop:path_based}
    The rank condition in Theorem~\ref{thm:minimal_excitation} is generically satisfied if there exist $|\bar B|$ vertex-disjoint paths in the network graph, each starting from a node in $r^A$ and $\bar r^B$ and ending at a \textbf{distinct} output in $\bar y^B$.
\end{proposition}

Similar propositions derived from the main result in \cite{graph_TF_rank} can be found in the literature, e.g. \cite[Proposition V.1]{HendrickGeversBazanella2019} and \cite[Proposition 1]{VandenHofRamaswamy2020}

Theorem~\ref{thm:minimal_excitation} together with Proposition~\ref{prop:path_based} show that complete excitation of all external signals in networks $B$ is generally not necessary. Instead, a carefully selected subset $\bar r^B$ suffices, with the selection guided by the network topology.

\bibliographystyle{plain}        % Include this if you use bibtex 
\bibliography{refs} 


\end{document}